\section{Methodology}

We describe our approach to measuring reward information bandwidth effects on code LLM training.

\subsection{Information Bandwidth Framework}

We frame reward granularity in information-theoretic terms (approximate bit estimates):

\textbf{LOW (Binary):} $r \in \{0, 1\}$ provides $\sim$1 bit per update.

\textbf{MEDIUM (Continuous):} $r = n_{\text{passed}} / n_{\text{total}}$ provides additional information via pass rate granularity.

\textbf{HIGH (Categorical + Continuous):} $r = 0.5 \times \text{categorical} + 0.3 \times \text{pass\_ratio} + 0.2 \times \text{partial\_credit}$ combines multiple signals.

\emph{Note: Precise bit calculations would require entropy analysis of the actual reward distributions, which we did not perform. The ``$\sim$1.5 bits'' and ``$\sim$2 bits'' estimates are heuristic.}

\subsection{Error-Type Scoring}

\begin{table}[h]
\centering
\begin{tabular}{@{}lll@{}}
\toprule
Error Type & Score & Rationale \\
\midrule
Syntax error & 0.00 & Code does not parse \\
Runtime error & 0.25 & Code runs but crashes \\
Assertion failure & 0.50 & Code runs but wrong output \\
Pass & 1.00 & Code is correct \\
\bottomrule
\end{tabular}
\caption{Error-type scoring scheme for categorical reward component.}
\label{tab:error-scoring}
\end{table}

\subsection{Training Configuration}

\begin{itemize}
    \item \textbf{Model:} CodeLlama-7B-Instruct
    \item \textbf{Training data:} MBPP sanitized ($\sim$374 problems)
    \item \textbf{Evaluation:} MBPP validation (50 problems)
    \item \textbf{Algorithm:} REINFORCE with advantage baseline (not PPO)
    \item \textbf{Epochs:} 1 (PoC scale)
    \item \textbf{Seeds:} 3 per condition
    \item \textbf{Batch size:} 4
    \item \textbf{Gradient steps:} $\sim$94 per epoch
\end{itemize}
